% =============================================================================
% Appendix A — ODD protocol (Grimm et al. 2020, second update).
% The anti-black-box deliverable: sufficient for reimplementation.
% =============================================================================
\section{ODD Protocol for the Cohort Microsimulation}
\label{app:odd}

This appendix documents the model of Section~\ref{sec:model} following the ODD
(Overview, Design concepts, Details) protocol, second update
\citep{grimm2020odd}. Together with Algorithm~\ref{alg:abm} and
equations~\eqref{eq:formation}--\eqref{eq:phi}, it is intended to be sufficient
for independent reimplementation. The reference implementation (Julia,
Agents.jl) with all figure data is available in the project repository.

\subsection{Purpose and patterns}

\emph{Purpose:} to test whether cohort-entry heterogeneity in marriage
propensity --- with no social feedback of any kind --- is \emph{sufficient} to
reproduce the observed union-composition collapse among women 20--39 in Costa
Rica (2010--2024) and Colombia (2008--2024). The model is evaluated against
three pre-registered patterns: (M1) the magnitude of the married-share fall;
(M2) its late-loading; (M3) the sign of its late-window curvature
(Section~\ref{sec:calibration}). It is explicitly \emph{not} used for
identification (done econometrically, Section~\ref{sec:discrimination}),
fertility-magnitude claims, or projection.

\subsection{Entities, state variables, and scales}

\emph{Entities:} women; no other agent types, no spatial units, no collectives.
\emph{State variables per woman:} age $a \in \{15,\dots,49\}$ (single years);
birth year $b$ (fixed at entry); union status
$u \in \{\Sng,\Coh,\Mar\}$; parity (bookkeeping for the comparison-only
fertility overlay; no behavioral role). Two fixed background attributes
(education level, urban/rural) are carried for descriptive seeding realism but
enter no transition rule in this specification.
\emph{Scales:} one tick = one calendar year; population 50{,}000 women
(reported ensembles); horizon = seed year to 2024 (14 ticks Costa Rica, 16
Colombia).

\subsection{Process overview and scheduling}

Each tick, every woman executes, in order: (1) union transition per
equations~\eqref{eq:formation}--\eqref{eq:dissolution}; (2) fertility-overlay
draw (comparison bookkeeping only); (3) ageing by one year, with closed-cohort
replacement --- a woman exceeding 49 is replaced by a new 15-year-old single
woman with birth year $t-15$, who therefore carries her own cohort's $\kb$ and
$\phib$. The model-level step advances the calendar year; it computes nothing
else. Agent processing order is immaterial because agents do not interact
(Design concepts). Composition observables are measured after each tick.

\subsection{Design concepts}

\emph{Basic principles.} Cohort replacement as the carrier of behavioral
change: transition probabilities depend on own birth cohort and own state only.
The two cohort functions $\kb$ (marriage entry) and $\phib$ (union formation)
encode the Second Demographic Transition's across-cohort gradient in reduced
form.
\emph{Emergence.} The aggregate composition trajectory --- including its late
acceleration --- emerges from the interaction of the cohort gradient with the
age-structured stock-flow dynamics: nothing in the rules accelerates; the
acceleration is compositional.
\emph{Adaptation, objectives, learning, prediction.} None. Agents do not
decide, optimize, learn, or forecast; transitions are exogenous hazards. This
is by design --- the model is a sufficiency probe for a mechanism, not a
behavioral theory.
\emph{Sensing and interaction.} \textbf{None, verifiably.} No agent's
transition depends on any other agent's state or on any population aggregate.
The absence is enforced by automated checks: static source checks (the agent
update contains no cross-agent read; feedback identifiers from the rejected
threshold model are absent) and a runtime registry proving that only the
composition target file is opened during calibration. Eleven checks per
country; all pass.
\emph{Stochasticity.} Bernoulli draws for all transitions; seeded generators;
the full pipeline is deterministic given seeds (the Costa Rican calibration
reproduces its loss to all reported digits on re-run).
\emph{Collectives.} None.
\emph{Observation.} Per-band (20--24, 25--29, 30--34, 35--39) married and
cohabiting shares each tick; ensemble mean $\pm$ s.d.\ over 16 seeds.

\subsection{Initialization}

The population is seeded at the first observed year (Costa Rica 2010, Colombia
2008 --- the 2007 Colombian survey year is a documented frame outlier,
Section~\ref{sec:datalim-measurement}): ages uniform 15--49; birth year $=$
seed year $-$ age; union status drawn from the observed band-specific
married/cohabiting shares, with age-graded extrapolation to the unobserved
15--19 and 40--49 tails (ramp below the 20--24 band; held at the 35--39 level,
respectively). Pre-seed cohorts receive $\kb$ and $\phib$ from the same
parametric functions --- the seeded stock embodies the gradient.

\subsection{Input data}

Calibration target (the only file read during calibration): observed annual
per-band married/cohabiting shares --- Costa Rica ENAHO via INEC REDATAM
tabulations; Colombia DANE GEIH person microdata, pooled annual files,
person-weighted. National TFR series are loaded only after calibration, for
the comparison overlay. Every series carries per-row source and methodology
flags in the repository.

\subsection{Submodels and calibrated parameters}

The complete submodels are
equations~\eqref{eq:formation}--\eqref{eq:phi}. Table~\ref{tab:params} reports
the calibrated headline (cohort-only) parameters; the robustness specification
adds the period shifter $\pit$ (three parameters), whose ablation costs are
reported in Section~\ref{sec:credibility}.

\begin{table}[h]
\centering
\caption{Calibrated parameters, headline (cohort-only) specification.}
\label{tab:params}
\small
\begin{tabular}{llcc}
\toprule
& Interpretation & Costa Rica & Colombia \\
\midrule
$f_0$ & base union-formation hazard (single, /yr) & 0.067 & 0.126 \\
$m_0$ & married share of new unions (oldest cohorts) & 0.308 & 0.147 \\
$c_0$ & cohabitation$\to$marriage hazard (oldest cohorts) & 0.042 & 0.025 \\
$\delta_c$ & cohabitation dissolution hazard & 0.062 & 0.067 \\
$\delta_m$ & marriage dissolution hazard & 0.019 & 0.023 \\
\midrule
$\kappa_\infty$ & marriage-propensity floor, newest cohorts & 0.371 & 0.306 \\
$b_0$ & cohort-decline midpoint (birth year) & 1991.7 & 1988.8 \\
$s$ & cohort-decline steepness (/birth-yr) & 0.129 & 0.194 \\
$\tau$ & formation tilt (/birth-yr) & 0.028 & 0.020 \\
\midrule
& composition loss~\eqref{eq:loss} & 0.0372 & 0.0246 \\
\bottomrule
\end{tabular}
\end{table}

\subsection{Data and code availability}

The reference implementation, calibration harness (including the
pre-registered gates and the automated no-feedback/no-TFR checks with their
logs), figure data for every exhibit, and the archived superseded runs are
maintained in the project repository and will be deposited with the
publication version. Figures regenerate from stored figure-data files without
re-simulation; the full calibration re-runs in minutes on a workstation and is
deterministic given seeds.
